% Member 3 report sections: supply (HLD 5 services 7, 8, 9, 6 and 2; HLD 6.2 Replenishment; HLD 6.3 Inbound; the
% forecast model comparison of task 3.3 and the route lengths of task 3.6; one HLD 13 trade-off).
% Body only, so the team report can % Member 3 report sections: supply (HLD 5 services 7, 8, 9, 6 and 2; HLD 6.2 Replenishment; HLD 6.3 Inbound; the
% forecast model comparison of task 3.3 and the route lengths of task 3.6; one HLD 13 trade-off).
% Body only, so the team report can % Member 3 report sections: supply (HLD 5 services 7, 8, 9, 6 and 2; HLD 6.2 Replenishment; HLD 6.3 Inbound; the
% forecast model comparison of task 3.3 and the route lengths of task 3.6; one HLD 13 trade-off).
% Body only, so the team report can % Member 3 report sections: supply (HLD 5 services 7, 8, 9, 6 and 2; HLD 6.2 Replenishment; HLD 6.3 Inbound; the
% forecast model comparison of task 3.3 and the route lengths of task 3.6; one HLD 13 trade-off).
% Body only, so the team report can \input{member3-supply}. Needs: array, booktabs, tikz (arrows.meta, positioning,
% calc).
% member3-standalone.tex compiles these sections on their own.
%
% Numbers measured on main (8c83ed9) on Sunday 22 November 2026, on throwaway SQLite databases. The generated sales
% are drawn back from the day the script runs, so another day gives slightly different numbers. To re-run them, from
% the repository root with the project's virtual environment:
%   export DARKSTORE_DB=sqlite:///$(mktemp -d)/m3.db BACKGROUND_JOBS=off
%   python -m app.seed && python scripts/generate_sales.py            # demo data and 60 days of sales (seed 42)
%   python -c "from app.db import SessionLocal; from app.modules.forecasting.service import compare_models as c; db = SessionLocal(); [print(s, d, c(db, s, d)['wape']) for s in (1, 2) for d in (14, 28)]"
%     -> Table tab:m3-forecast. For the ranges in the text, repeat the three lines with generate_sales.py --seed 1 to 10.
%   python -c "from app.db import SessionLocal; from app.modules.forecasting import service as f; from app.modules.replenishment import service as r; db = SessionLocal(); print([x for x in f.run(db, 1) if x['sku_id'] == 1]); print(r.order_view(r.run(db, 1)['transfer_order'])); print(r.run(db, 1))"
%     -> the milk example in section sec:m3-replenishment (after the comparison: it changes the reorder points).
%   python scripts/compare_slotting.py --orders 200 --seed 7            # uses its own temporary database
%   python scripts/compare_slotting.py --orders 5000 --seed 7           # -> Table tab:m3-slotting; --seed 1 to 10 for the range

\section{Supply services (HLD 5, services 7, 8, 9, 6 and 2)}
\label{sec:m3-services}

Members 1 and 2 make sure that an order for something in stock goes well. Supply decides whether it is in stock in
the first place: the shelf should be full before the customer asks. Four services form a loop
(Figure~\ref{fig:m3-loop}), and Assortment sits above it, deciding which SKUs a store carries at all. None of the five
writes a stock quantity: Inbound asks Inventory to receive and Slotting asks it to move, so Member 1's three stock
records always agree. Forecasting writes only two planning settings on each stock row, the reorder point and the
target stock.

\begin{figure}[ht]
  \centering
  \begin{tikzpicture}[
      >={Stealth[length=2.2mm]}, node distance=17mm,
      svc/.style={draw=black!45, rounded corners=3pt, minimum width=25mm, minimum height=11mm, align=center,
                  font=\bfseries\small},
      shelf/.style={svc, draw=blue!70!black, fill=blue!8, line width=0.9pt, minimum width=62mm},
      edgelabel/.style={font=\footnotesize, text=black!65, align=center}]
    \node[svc] (fc) {Forecasting\\[-1pt]{\mdseries\scriptsize service 7}};
    \node[svc, right=of fc] (rp) {Replenishment\\[-1pt]{\mdseries\scriptsize service 8}};
    \node[svc, right=of rp] (ib) {Inbound\\[-1pt]{\mdseries\scriptsize service 9}};
    \node[svc, right=of ib] (sl) {Slotting\\[-1pt]{\mdseries\scriptsize service 6}};
    \node[svc, above=10mm of rp] (as) {Assortment\\[-1pt]{\mdseries\scriptsize service 2}};
    \node[shelf, below=16mm of $(rp)!0.5!(ib)$] (shelf) {Stock in bins\\[-1pt]{\mdseries\scriptsize Inventory
      (Member 1) writes every change}};

    \draw[->, black!55] (fc) -- node[edgelabel, above] {reorder\\point} (rp);
    \draw[->, black!55] (rp) -- node[edgelabel, above] {transfer\\order} (ib);
    \draw[->, black!55] (sl) -- node[edgelabel, above] {bin per\\line} (ib);
    \draw[->, black!55] (as) -- node[edgelabel, right] {SKUs carried} (rp);
    \draw[->, black!55] (ib.south) -- node[edgelabel, right] {put away} (ib.south |- shelf.north);
    \draw[->, black!55] (sl.south) |- node[edgelabel, pos=0.3, right] {moves} (shelf.east);
    \draw[->, black!55] (shelf.west) -| node[edgelabel, pos=0.75, left] {sales\\(ledger)} (fc.south);
  \end{tikzpicture}
  \caption{The supply loop. Only Inventory changes stock: Inbound and Slotting ask it to receive and to move units.}
  \label{fig:m3-loop}
\end{figure}

\subsection{Forecasting (service 7)}

Forecasting predicts each SKU's daily sales and turns them into a reorder point (FR-7). Sales come from the
\texttt{sale} rows of the ledger, in whole days ending yesterday, because today is only partly over.

\begin{itemize}
  \item \textbf{Three models with one shape.} Each forecasts a day from the days before it: the \emph{moving average}
    of the last 7 days, the \emph{weekday average} of the same weekday in the last 4 weeks, and \emph{exponential
    smoothing} ($\alpha = 0.3$), in which each earlier day counts 0.7 times as much as the day after it. Days before a
    SKU's first sale are unknown, not zero, so a new SKU is not forecast as if it had sold nothing for weeks.
  \item \textbf{A measuring stick first.} The backtest replays the last 14 days, forecasting each day from the days
    before it, and reports the error over every SKU and day,
    \[ \mathrm{WAPE} = \frac{\sum |\textrm{actual} - \textrm{forecast}|}{\sum \textrm{actual}}, \]
    for which HLD 12 asks 30\% or lower. Weighting by volume stops slow sellers swamping the score. A day counts only
    when the SKU has a week of sales before it, so on the demo data's single week of sales the answer is ``not
    available'', not a misleading number. A comparison runs all three models on the same days, and the setting
    \texttt{FORECAST\_MODEL} picks the one in use: the weekday average, which won (section~\ref{sec:m3-results}).
  \item \textbf{Safety stock from the model's errors.} The reorder point uses the spread of the model's \emph{errors}
    over the last 14 days, not the spread of raw sales: safety stock should cover what the forecast gets wrong, not
    variation it already predicts, such as busier weekends.
\end{itemize}
A manager runs it for a store (\texttt{POST /v1/forecasting/run}), and each forecast is saved with its model's name.

\subsection{Replenishment (service 8)}

Replenishment orders from the mother hub when stock runs low (FR-8), in whole cases and within each zone's free bin
space. It runs every 15 minutes for every open store, and whenever a manager asks; section~\ref{sec:m3-replenishment}
follows one run. It also publishes \texttt{stock.low} when the stock not held for customers falls below the reorder point, lets
a manager stop reordering one SKU so that it runs down before it is dropped, and reports the five supply KPIs of HLD 12
(stock-outs, days of cover, wastage, deliveries on time and forecast WAPE) against their targets for Member 4's ops
dashboard. The hub itself is a stand-in: a button plays its warehouse system saying ``the truck has left''.

\subsection{Inbound (service 9)}

Inbound receives a truck in the five steps of HLD 6.3 (section~\ref{sec:m3-inbound}). Each truck gets a saved goods
receipt whose lines record, per SKU and batch, the units scanned, damaged, refused and put away. Putting a line away
calls \texttt{inventory.receive} (Member 1), which writes the bin, the totals row and the ledger in one transaction.

\subsection{Slotting (service 6)}

Slotting decides which bin a SKU lives in. Each bin has a zone (ambient, chilled or frozen), a capacity and a
\texttt{walk\_seq}, its place on the picker's route, where bin 1 is next to the packing desk. HLD service 6 ``places
fast movers near the packing station'':
\begin{itemize}
  \item \textbf{The plan} (look only) ranks each zone's SKUs by units sold a day over the last 14 days and gives the
    fastest the bin nearest the desk, the next one the next bin, and so on. A SKU that sold nothing stays where it is
    unless a seller needs its bin, since moving it would not shorten any walk.
  \item \textbf{Apply} (managers) locks the store's bins and moves each batch through \texttt{inventory.move}: two
    ledger rows, nothing gained or lost. A batch moves only when its new bin has room at that moment, so no bin is ever
    over capacity, and a batch that a picker has been sent to stays where it is.
  \item \textbf{New deliveries} go to the SKU's planned bin first, then to bins already holding it, then to the
    emptiest bins of its zone, never past capacity, so the shelf re-sorts itself as stock arrives.
\end{itemize}
A SKU never leaves its own temperature zone; section~\ref{sec:m3-results} shows what that costs in walking.

\subsection{Assortment (service 2)}

Assortment decides which SKUs each store carries. Adding a SKU creates its stock row with zero stock, so
Replenishment starts ordering it; dropping one is refused while any unit is on the shelf, held for a customer or on its
way. From the last 30 days of sales it suggests SKUs to \emph{add} (selling in other stores, highest demand first),
SKUs to \emph{drop} (stock on hand 30 days ago and not one unit sold since) and the \emph{space} left in each zone
(bins against SKUs carried, so nobody adds a fifth chilled SKU to four chilled bins). A manager decides.

\section{Replenishment (HLD 6.2)}
\label{sec:m3-replenishment}

HLD 6.2 keeps every SKU above a reorder point. The HLD counts in hours and the prototype in days, with a lead time $L$
of 1 day, $z = 1.65$ (stock should run out in at most 5\% of cycles) and 3 days of cover:
\[
  \textrm{reorder point} = \bigl\lceil dL + z\sigma\sqrt{L}\, \bigr\rceil \qquad\qquad
  \textrm{target stock} = \textrm{reorder point} + \lceil 3d \rceil
\]
Here $d$ is the forecast daily demand and $\sigma$ the spread of the forecast's errors. The first part of the reorder
point covers the sales expected while the truck is on its way; the second is safety stock for a busier day.

\begin{enumerate}
  \item \textbf{Check.} Each SKU's \emph{position}, on hand $-$ held $+$ in transit, is compared with its reorder
    point. Counting what is already on its way is what makes the 15-minute timer harmless: without it, every run would
    order the same items again until the truck came.
  \item \textbf{Size.} A SKU below its reorder point needs target stock $-$ position, rounded up to whole cases
    (12 milk, 6 eggs, 10 bread).
  \item \textbf{Fit.} A zone's free space is its bins' capacity minus the units stored and the units on their way.
    When the lines do not fit, the most urgent go first: lowest days of cover (what customers can buy $\div$ daily
    demand), then the faster seller. The last line is cut to the cases that fit, and every cut is reported, also when
    the timer made it.
  \item \textbf{Order.} One transfer order takes every line (created, dispatched, received). Only one run happens at a
    time in an app process, so the timer and a manager's button cannot both order the same stock.
\end{enumerate}

\paragraph{One run, with real numbers.} On the generated sales of section~\ref{sec:m3-results}, run on Sunday
22 November, milk in store 1 is forecast at 32.25 a day: the mean of the last four Sundays, one of them a festival day.
The spread of its errors is 6.5, so the reorder point is $\lceil 32.25 + 1.65 \times 6.5 \rceil = 43$ and the target
$43 + \lceil 3 \times 32.25 \rceil = 140$. With 30 on hand, none held and none on the way, the run orders
$140 - 30 = 110$, rounded up to 10 cases of 12: 120. The same transfer order took bread (4 left: 9 cases of 10), eggs,
curd, paneer and ice cream, and nothing was cut: the three chilled lines needed 362 of the 510 free chilled units. A
second run straight after ordered nothing.

The example also shows a weakness: the reorder point is sized from \emph{today's} forecast only. Run on a Sunday after
a festival Sunday, milk is sized for 32 a day; run on a weekday, it would be sized for a weekday.

\section{Inbound (HLD 6.3)}
\label{sec:m3-inbound}

HLD 6.3 receives a truck in five steps and asks that ``each item becomes sellable as soon as its bin scan is
confirmed, not when the whole truck is done''. In the prototype only step 4 changes stock.

\begin{enumerate}
  \item \textbf{Open.} A receipt is saved for the transfer order. The order is locked first, so a double click on
    \emph{Receive truck} opens one receipt, not two.
  \item \textbf{Scan.} Each case is scanned with its batch number and expiry date, which first-expiry-first-out
    picking relies on. Good units wait on the dock, not yet for sale. Damaged units, and a batch with less than half
    its shelf life left (milk, which keeps 3 days, with 1 day to go), are recorded but never stocked.
  \item \textbf{Putaway tasks.} Slotting names a bin, or several, for each line on the dock, never past capacity.
  \item \textbf{Put away.} The units go into the bin through \texttt{inventory.receive} and are \emph{for sale as soon
    as this commits}, whatever is still on the dock. They stop counting as in transit at the same moment.
  \item \textbf{Close.} The transfer order is marked received, and every SKU is reported as ordered against received,
    with damaged, refused and sent-back units. Closing is refused while good units wait on the dock, unless the
    receiver sends them back to the hub.
\end{enumerate}

A receipt opened by mistake can be cancelled until anything from it is put away, and the real truck can still be
received. The tests check the property HLD 6.3 asks for: after a scan, what customers can buy is unchanged; after one
line is put away, that SKU rises and no other does.

\section{Measurements}
\label{sec:m3-results}

Both measurements run the project's own code on a throwaway copy of the demo data with 60 days of generated sales
(\texttt{scripts/generate\_sales.py}, task 3.1, seed 42): for every SKU and day, a random (Poisson) number of units
around a base demand per category (20 a day for dairy, including eggs, 12 for bread, 4 for staples, 2 for frozen
food, 3 for the rest), 35\% more at weekends and twice as much on two festival days. They were measured on main on
22 November 2026; the commands are at the top of \texttt{member3-supply.tex}. The history is drawn back from the day
it is generated, so another day gives slightly different numbers.

\subsection{Forecast models (task 3.3)}

\begin{table}[ht]
  \centering
  \footnotesize
  \begin{tabular}{@{}lrrrr@{}}
    \toprule
    Model & \shortstack[r]{Store 1\\14 days} & \shortstack[r]{Store 2\\14 days} & \shortstack[r]{Store 1\\28 days}
          & \shortstack[r]{Store 2\\28 days} \\
    \midrule
    Moving average, last 7 days (the baseline) & 34.4\% & 36.1\% & 30.7\% & 32.4\% \\
    Weekday average, last 4 same weekdays      & \textbf{29.3\%} & \textbf{29.2\%} & \textbf{28.9\%} & \textbf{27.8\%} \\
    Exponential smoothing, $\alpha = 0.3$      & 35.9\% & 39.2\% & 32.4\% & 34.3\% \\
    \bottomrule
  \end{tabular}
  \caption{Forecast error (WAPE, lower is better) of each model over the last 14 and 28 days, every SKU of the store
    (\texttt{GET /v1/forecasting/compare}). HLD 12 target: 30\% or lower.}
  \label{tab:m3-forecast}
\end{table}

The weekday average wins in both stores and both windows (Table~\ref{tab:m3-forecast}). Over 14 days it is 5 to 7
points better than the baseline and the only model within the 30\% target, so it is the model in use. With ten other
random seeds for the history (40 more store-and-window cases) it was again the best every time, at 26.0 to 32.3\%
against 30.5 to 39.1\% for the moving average and 32.1 to 41.1\% for smoothing, and it met the target in 28 of the 40.

\begin{itemize}
  \item \textbf{Why it wins.} A 7-day moving average forecasts every day as the week's average, so with weekends 35\%
    busier it is about 10\% too high on weekdays and 19\% too low at weekends. The weekday average forecasts a
    Saturday from Saturdays. Smoothing does worst, because it leans on the last two or three days: a Monday is
    forecast mostly from the weekend, and a Saturday from the weekdays.
  \item \textbf{What is left is mostly noise.} For a SKU that sells 3 a day, one unit off is already a third of the
    day's sales, so the slow sellers keep the store near 30\% whatever the model. The 14-day window also starts on a
    festival day, which every model misses; the weekday average then carries it into the same weekday's forecast for
    four weeks, as in the milk example of section~\ref{sec:m3-replenishment}.
\end{itemize}

\subsection{Velocity slotting (task 3.6)}

\texttt{scripts/compare\_slotting.py} draws random orders of 1 to 3 different items, each item as likely as its share
of the last 14 days' sales, and finds the bin each item would be picked from, as Inventory sends a picker: first expiry
first, then nearest. A walk costs twice the furthest \texttt{walk\_seq} it visits (Member 2's model, task 2.6). The
script then applies the plan and measures the same orders again (Table~\ref{tab:m3-slotting}).

\begin{table}[ht]
  \centering
  \footnotesize
  \begin{tabular}{@{}llrrrr@{}}
    \toprule
    Run & Orders measured & Orders & Before & After & Saving \\
    \midrule
    200 orders (seed 7)     & All orders         & 200     & 24.6 & 24.3 & 1\%  \\
                            & Ambient items only & 49      & 11.2 &  9.6 & 15\% \\
    \addlinespace[3pt]
    5{,}000 orders (seed 7) & All orders         & 5{,}000 & 24.1 & 23.9 & 1\%  \\
                            & Ambient items only & 1{,}317 & 10.6 &  9.9 & 6\%  \\
    \bottomrule
  \end{tabular}
  \caption{Average walk per order in bin steps, store 1, before and after the re-slotting plan was applied (13 batches
    moved, none skipped). The guide asks for 200 orders; 5,000 are less noisy.}
  \label{tab:m3-slotting}
\end{table}

The plan shortens the average walk by only 1\%. Orders of ambient items alone save 6\%; the 15\% of the 200-order run
is luck, and order seeds 1 to 10 gave between $-1$\% and 2\% for all orders and between $-9$\% and 15\% for ambient
ones. The plan for store 1 shows why:
\begin{itemize}
  \item \textbf{The fridges decide most walks.} Slotting re-sorts bins only within a temperature zone, and the chilled
    and frozen bins are at the far end of the walk (\texttt{walk\_seq} 13 to 18). Milk, curd and paneer sell 22 to 24
    a day each, 47\% of all units sold, so about three orders in four (3{,}683 of 5{,}000) include a chilled or frozen
    item and walk to the fridges whatever the plan.
  \item \textbf{The fastest ambient items were already at the front.} The demo layout had eggs (24.9 a day) and bread
    (14.8) in the first two bins; the plan only swaps them.
  \item \textbf{Equal sellers only trade places.} The three chilled items sell at almost the same rate, so curd and
    milk only swap the first two chilled bins. The other ambient items sell 3 to 5 a day, so their order is mostly noise.
\end{itemize}
In this layout the walk is decided by where the fridges stand, not by which shelf holds which item. A bigger saving
needs the fridges nearer the packing desk, which is a layout decision. Slotting should matter more in a real store,
with thousands of SKUs of very different speeds in each zone, than in a demo store with 16.

\section{Trade-off: for sale at putaway, not when the truck is done (HLD 13)}
\label{sec:m3-tradeoff}

The prototype first received a truck in one call: one transaction, all or nothing, and no record of the receipt. That
is the simplest correct design, because if anything fails, nothing happened. HLD 6.3 asks for more, each line for sale
the moment it is in its bin, and task 3.5 built that as five steps with a saved receipt. Building it showed what a
faster shelf costs.

\textbf{What it gains.} A fast seller that has run out is for sale as soon as its case is on the shelf, not after the
last case of the truck is scanned. The receipt also records what was ordered, scanned, damaged, refused, put away and
sent back, which the single call never kept.

\textbf{What it costs.}
\begin{itemize}
  \item \textbf{A third place for stock.} Units on the dock are neither on the shelf nor gone, and every rule had to say
    what happens to them: they still count as in transit, a receipt cannot close while they wait unless they are sent
    back, and a receipt can be cancelled only before anything from it is put away.
  \item \textbf{Races and retries the single call never had.} Two receivers could put away the same line, one could
    close the receipt while another scans, and a handheld could send a scan twice. Each step locks the receipt's row,
    units leave the dock through one conditional \texttt{UPDATE}, a bin's room is checked again after writing, and
    every write takes an \texttt{Idempotency-Key}, so a repeated scan counts once.
  \item \textbf{Edge cases keep appearing.} Review found that a line whose batch number is already in the store with
    another expiry date can be scanned onto the dock but never put away, so its receipt closes only by sending the line
    back. The fix, refusing such a scan at the dock, is waiting for review.
\end{itemize}

We kept the single call as well: \texttt{POST /v1/goods-receipts} runs the five steps in one transaction, so the ops
screen's \emph{Receive truck} button is still all or nothing, and both paths share one set of rules. We accepted a receipt with more states, and more
code to keep them consistent, in exchange for a shelf that refills line by line and a record of every truck.

\section{Limits and open items}

Table~\ref{tab:m3-limits} lists what the supply work leaves open, and what would close each item.

\begin{table}[!htb]
  \centering
  \footnotesize
  \begin{tabular}{@{}>{\raggedright\arraybackslash}p{0.27\linewidth}>{\raggedright\arraybackslash}p{0.36\linewidth}
                  >{\raggedright\arraybackslash}p{0.31\linewidth}@{}}
    \toprule
    Limit & What it means & Next step \\
    \midrule
    Daily forecasts, not hourly & The evening peak is not planned for & Spread each daily forecast over the hours by
      their share of sales \\
    \addlinespace[3pt]
    No festivals or weather & A festival day is missed, then repeated in the same weekday's forecast for four weeks
      & A festival calendar as an input \\
    \addlinespace[3pt]
    Reorder point from today's forecast only & Run on a busy Sunday, milk is sized for a Sunday & Size it from the
      forecasts of the lead-time days \\
    \addlinespace[3pt]
    The same cover for fresh and long-life food & Three days of cover above the reorder point, also for milk that keeps
      three days & Cover days that depend on shelf life \\
    \addlinespace[3pt]
    A batch number already in the store with another expiry date & It is scanned onto the dock but cannot be put away;
      only sending it back closes the receipt & Refuse it at the scan (pull request \#77, pending) \\
    \addlinespace[3pt]
    Slotting only within a zone; apply records the moves at once & 1\% shorter walks in this layout, and staff must
      carry the stock straight away & Fridges nearer the packing desk; items bought together side by side \\
    \bottomrule
  \end{tabular}
  \caption{Limits of the supply work and what would close them.}
  \label{tab:m3-limits}
\end{table}
. Needs: array, booktabs, tikz (arrows.meta, positioning,
% calc).
% member3-standalone.tex compiles these sections on their own.
%
% Numbers measured on main (8c83ed9) on Sunday 22 November 2026, on throwaway SQLite databases. The generated sales
% are drawn back from the day the script runs, so another day gives slightly different numbers. To re-run them, from
% the repository root with the project's virtual environment:
%   export DARKSTORE_DB=sqlite:///$(mktemp -d)/m3.db BACKGROUND_JOBS=off
%   python -m app.seed && python scripts/generate_sales.py            # demo data and 60 days of sales (seed 42)
%   python -c "from app.db import SessionLocal; from app.modules.forecasting.service import compare_models as c; db = SessionLocal(); [print(s, d, c(db, s, d)['wape']) for s in (1, 2) for d in (14, 28)]"
%     -> Table tab:m3-forecast. For the ranges in the text, repeat the three lines with generate_sales.py --seed 1 to 10.
%   python -c "from app.db import SessionLocal; from app.modules.forecasting import service as f; from app.modules.replenishment import service as r; db = SessionLocal(); print([x for x in f.run(db, 1) if x['sku_id'] == 1]); print(r.order_view(r.run(db, 1)['transfer_order'])); print(r.run(db, 1))"
%     -> the milk example in section sec:m3-replenishment (after the comparison: it changes the reorder points).
%   python scripts/compare_slotting.py --orders 200 --seed 7            # uses its own temporary database
%   python scripts/compare_slotting.py --orders 5000 --seed 7           # -> Table tab:m3-slotting; --seed 1 to 10 for the range

\section{Supply services (HLD 5, services 7, 8, 9, 6 and 2)}
\label{sec:m3-services}

Members 1 and 2 make sure that an order for something in stock goes well. Supply decides whether it is in stock in
the first place: the shelf should be full before the customer asks. Four services form a loop
(Figure~\ref{fig:m3-loop}), and Assortment sits above it, deciding which SKUs a store carries at all. None of the five
writes a stock quantity: Inbound asks Inventory to receive and Slotting asks it to move, so Member 1's three stock
records always agree. Forecasting writes only two planning settings on each stock row, the reorder point and the
target stock.

\begin{figure}[ht]
  \centering
  \begin{tikzpicture}[
      >={Stealth[length=2.2mm]}, node distance=17mm,
      svc/.style={draw=black!45, rounded corners=3pt, minimum width=25mm, minimum height=11mm, align=center,
                  font=\bfseries\small},
      shelf/.style={svc, draw=blue!70!black, fill=blue!8, line width=0.9pt, minimum width=62mm},
      edgelabel/.style={font=\footnotesize, text=black!65, align=center}]
    \node[svc] (fc) {Forecasting\\[-1pt]{\mdseries\scriptsize service 7}};
    \node[svc, right=of fc] (rp) {Replenishment\\[-1pt]{\mdseries\scriptsize service 8}};
    \node[svc, right=of rp] (ib) {Inbound\\[-1pt]{\mdseries\scriptsize service 9}};
    \node[svc, right=of ib] (sl) {Slotting\\[-1pt]{\mdseries\scriptsize service 6}};
    \node[svc, above=10mm of rp] (as) {Assortment\\[-1pt]{\mdseries\scriptsize service 2}};
    \node[shelf, below=16mm of $(rp)!0.5!(ib)$] (shelf) {Stock in bins\\[-1pt]{\mdseries\scriptsize Inventory
      (Member 1) writes every change}};

    \draw[->, black!55] (fc) -- node[edgelabel, above] {reorder\\point} (rp);
    \draw[->, black!55] (rp) -- node[edgelabel, above] {transfer\\order} (ib);
    \draw[->, black!55] (sl) -- node[edgelabel, above] {bin per\\line} (ib);
    \draw[->, black!55] (as) -- node[edgelabel, right] {SKUs carried} (rp);
    \draw[->, black!55] (ib.south) -- node[edgelabel, right] {put away} (ib.south |- shelf.north);
    \draw[->, black!55] (sl.south) |- node[edgelabel, pos=0.3, right] {moves} (shelf.east);
    \draw[->, black!55] (shelf.west) -| node[edgelabel, pos=0.75, left] {sales\\(ledger)} (fc.south);
  \end{tikzpicture}
  \caption{The supply loop. Only Inventory changes stock: Inbound and Slotting ask it to receive and to move units.}
  \label{fig:m3-loop}
\end{figure}

\subsection{Forecasting (service 7)}

Forecasting predicts each SKU's daily sales and turns them into a reorder point (FR-7). Sales come from the
\texttt{sale} rows of the ledger, in whole days ending yesterday, because today is only partly over.

\begin{itemize}
  \item \textbf{Three models with one shape.} Each forecasts a day from the days before it: the \emph{moving average}
    of the last 7 days, the \emph{weekday average} of the same weekday in the last 4 weeks, and \emph{exponential
    smoothing} ($\alpha = 0.3$), in which each earlier day counts 0.7 times as much as the day after it. Days before a
    SKU's first sale are unknown, not zero, so a new SKU is not forecast as if it had sold nothing for weeks.
  \item \textbf{A measuring stick first.} The backtest replays the last 14 days, forecasting each day from the days
    before it, and reports the error over every SKU and day,
    \[ \mathrm{WAPE} = \frac{\sum |\textrm{actual} - \textrm{forecast}|}{\sum \textrm{actual}}, \]
    for which HLD 12 asks 30\% or lower. Weighting by volume stops slow sellers swamping the score. A day counts only
    when the SKU has a week of sales before it, so on the demo data's single week of sales the answer is ``not
    available'', not a misleading number. A comparison runs all three models on the same days, and the setting
    \texttt{FORECAST\_MODEL} picks the one in use: the weekday average, which won (section~\ref{sec:m3-results}).
  \item \textbf{Safety stock from the model's errors.} The reorder point uses the spread of the model's \emph{errors}
    over the last 14 days, not the spread of raw sales: safety stock should cover what the forecast gets wrong, not
    variation it already predicts, such as busier weekends.
\end{itemize}
A manager runs it for a store (\texttt{POST /v1/forecasting/run}), and each forecast is saved with its model's name.

\subsection{Replenishment (service 8)}

Replenishment orders from the mother hub when stock runs low (FR-8), in whole cases and within each zone's free bin
space. It runs every 15 minutes for every open store, and whenever a manager asks; section~\ref{sec:m3-replenishment}
follows one run. It also publishes \texttt{stock.low} when the stock not held for customers falls below the reorder point, lets
a manager stop reordering one SKU so that it runs down before it is dropped, and reports the five supply KPIs of HLD 12
(stock-outs, days of cover, wastage, deliveries on time and forecast WAPE) against their targets for Member 4's ops
dashboard. The hub itself is a stand-in: a button plays its warehouse system saying ``the truck has left''.

\subsection{Inbound (service 9)}

Inbound receives a truck in the five steps of HLD 6.3 (section~\ref{sec:m3-inbound}). Each truck gets a saved goods
receipt whose lines record, per SKU and batch, the units scanned, damaged, refused and put away. Putting a line away
calls \texttt{inventory.receive} (Member 1), which writes the bin, the totals row and the ledger in one transaction.

\subsection{Slotting (service 6)}

Slotting decides which bin a SKU lives in. Each bin has a zone (ambient, chilled or frozen), a capacity and a
\texttt{walk\_seq}, its place on the picker's route, where bin 1 is next to the packing desk. HLD service 6 ``places
fast movers near the packing station'':
\begin{itemize}
  \item \textbf{The plan} (look only) ranks each zone's SKUs by units sold a day over the last 14 days and gives the
    fastest the bin nearest the desk, the next one the next bin, and so on. A SKU that sold nothing stays where it is
    unless a seller needs its bin, since moving it would not shorten any walk.
  \item \textbf{Apply} (managers) locks the store's bins and moves each batch through \texttt{inventory.move}: two
    ledger rows, nothing gained or lost. A batch moves only when its new bin has room at that moment, so no bin is ever
    over capacity, and a batch that a picker has been sent to stays where it is.
  \item \textbf{New deliveries} go to the SKU's planned bin first, then to bins already holding it, then to the
    emptiest bins of its zone, never past capacity, so the shelf re-sorts itself as stock arrives.
\end{itemize}
A SKU never leaves its own temperature zone; section~\ref{sec:m3-results} shows what that costs in walking.

\subsection{Assortment (service 2)}

Assortment decides which SKUs each store carries. Adding a SKU creates its stock row with zero stock, so
Replenishment starts ordering it; dropping one is refused while any unit is on the shelf, held for a customer or on its
way. From the last 30 days of sales it suggests SKUs to \emph{add} (selling in other stores, highest demand first),
SKUs to \emph{drop} (stock on hand 30 days ago and not one unit sold since) and the \emph{space} left in each zone
(bins against SKUs carried, so nobody adds a fifth chilled SKU to four chilled bins). A manager decides.

\section{Replenishment (HLD 6.2)}
\label{sec:m3-replenishment}

HLD 6.2 keeps every SKU above a reorder point. The HLD counts in hours and the prototype in days, with a lead time $L$
of 1 day, $z = 1.65$ (stock should run out in at most 5\% of cycles) and 3 days of cover:
\[
  \textrm{reorder point} = \bigl\lceil dL + z\sigma\sqrt{L}\, \bigr\rceil \qquad\qquad
  \textrm{target stock} = \textrm{reorder point} + \lceil 3d \rceil
\]
Here $d$ is the forecast daily demand and $\sigma$ the spread of the forecast's errors. The first part of the reorder
point covers the sales expected while the truck is on its way; the second is safety stock for a busier day.

\begin{enumerate}
  \item \textbf{Check.} Each SKU's \emph{position}, on hand $-$ held $+$ in transit, is compared with its reorder
    point. Counting what is already on its way is what makes the 15-minute timer harmless: without it, every run would
    order the same items again until the truck came.
  \item \textbf{Size.} A SKU below its reorder point needs target stock $-$ position, rounded up to whole cases
    (12 milk, 6 eggs, 10 bread).
  \item \textbf{Fit.} A zone's free space is its bins' capacity minus the units stored and the units on their way.
    When the lines do not fit, the most urgent go first: lowest days of cover (what customers can buy $\div$ daily
    demand), then the faster seller. The last line is cut to the cases that fit, and every cut is reported, also when
    the timer made it.
  \item \textbf{Order.} One transfer order takes every line (created, dispatched, received). Only one run happens at a
    time in an app process, so the timer and a manager's button cannot both order the same stock.
\end{enumerate}

\paragraph{One run, with real numbers.} On the generated sales of section~\ref{sec:m3-results}, run on Sunday
22 November, milk in store 1 is forecast at 32.25 a day: the mean of the last four Sundays, one of them a festival day.
The spread of its errors is 6.5, so the reorder point is $\lceil 32.25 + 1.65 \times 6.5 \rceil = 43$ and the target
$43 + \lceil 3 \times 32.25 \rceil = 140$. With 30 on hand, none held and none on the way, the run orders
$140 - 30 = 110$, rounded up to 10 cases of 12: 120. The same transfer order took bread (4 left: 9 cases of 10), eggs,
curd, paneer and ice cream, and nothing was cut: the three chilled lines needed 362 of the 510 free chilled units. A
second run straight after ordered nothing.

The example also shows a weakness: the reorder point is sized from \emph{today's} forecast only. Run on a Sunday after
a festival Sunday, milk is sized for 32 a day; run on a weekday, it would be sized for a weekday.

\section{Inbound (HLD 6.3)}
\label{sec:m3-inbound}

HLD 6.3 receives a truck in five steps and asks that ``each item becomes sellable as soon as its bin scan is
confirmed, not when the whole truck is done''. In the prototype only step 4 changes stock.

\begin{enumerate}
  \item \textbf{Open.} A receipt is saved for the transfer order. The order is locked first, so a double click on
    \emph{Receive truck} opens one receipt, not two.
  \item \textbf{Scan.} Each case is scanned with its batch number and expiry date, which first-expiry-first-out
    picking relies on. Good units wait on the dock, not yet for sale. Damaged units, and a batch with less than half
    its shelf life left (milk, which keeps 3 days, with 1 day to go), are recorded but never stocked.
  \item \textbf{Putaway tasks.} Slotting names a bin, or several, for each line on the dock, never past capacity.
  \item \textbf{Put away.} The units go into the bin through \texttt{inventory.receive} and are \emph{for sale as soon
    as this commits}, whatever is still on the dock. They stop counting as in transit at the same moment.
  \item \textbf{Close.} The transfer order is marked received, and every SKU is reported as ordered against received,
    with damaged, refused and sent-back units. Closing is refused while good units wait on the dock, unless the
    receiver sends them back to the hub.
\end{enumerate}

A receipt opened by mistake can be cancelled until anything from it is put away, and the real truck can still be
received. The tests check the property HLD 6.3 asks for: after a scan, what customers can buy is unchanged; after one
line is put away, that SKU rises and no other does.

\section{Measurements}
\label{sec:m3-results}

Both measurements run the project's own code on a throwaway copy of the demo data with 60 days of generated sales
(\texttt{scripts/generate\_sales.py}, task 3.1, seed 42): for every SKU and day, a random (Poisson) number of units
around a base demand per category (20 a day for dairy, including eggs, 12 for bread, 4 for staples, 2 for frozen
food, 3 for the rest), 35\% more at weekends and twice as much on two festival days. They were measured on main on
22 November 2026; the commands are at the top of \texttt{member3-supply.tex}. The history is drawn back from the day
it is generated, so another day gives slightly different numbers.

\subsection{Forecast models (task 3.3)}

\begin{table}[ht]
  \centering
  \footnotesize
  \begin{tabular}{@{}lrrrr@{}}
    \toprule
    Model & \shortstack[r]{Store 1\\14 days} & \shortstack[r]{Store 2\\14 days} & \shortstack[r]{Store 1\\28 days}
          & \shortstack[r]{Store 2\\28 days} \\
    \midrule
    Moving average, last 7 days (the baseline) & 34.4\% & 36.1\% & 30.7\% & 32.4\% \\
    Weekday average, last 4 same weekdays      & \textbf{29.3\%} & \textbf{29.2\%} & \textbf{28.9\%} & \textbf{27.8\%} \\
    Exponential smoothing, $\alpha = 0.3$      & 35.9\% & 39.2\% & 32.4\% & 34.3\% \\
    \bottomrule
  \end{tabular}
  \caption{Forecast error (WAPE, lower is better) of each model over the last 14 and 28 days, every SKU of the store
    (\texttt{GET /v1/forecasting/compare}). HLD 12 target: 30\% or lower.}
  \label{tab:m3-forecast}
\end{table}

The weekday average wins in both stores and both windows (Table~\ref{tab:m3-forecast}). Over 14 days it is 5 to 7
points better than the baseline and the only model within the 30\% target, so it is the model in use. With ten other
random seeds for the history (40 more store-and-window cases) it was again the best every time, at 26.0 to 32.3\%
against 30.5 to 39.1\% for the moving average and 32.1 to 41.1\% for smoothing, and it met the target in 28 of the 40.

\begin{itemize}
  \item \textbf{Why it wins.} A 7-day moving average forecasts every day as the week's average, so with weekends 35\%
    busier it is about 10\% too high on weekdays and 19\% too low at weekends. The weekday average forecasts a
    Saturday from Saturdays. Smoothing does worst, because it leans on the last two or three days: a Monday is
    forecast mostly from the weekend, and a Saturday from the weekdays.
  \item \textbf{What is left is mostly noise.} For a SKU that sells 3 a day, one unit off is already a third of the
    day's sales, so the slow sellers keep the store near 30\% whatever the model. The 14-day window also starts on a
    festival day, which every model misses; the weekday average then carries it into the same weekday's forecast for
    four weeks, as in the milk example of section~\ref{sec:m3-replenishment}.
\end{itemize}

\subsection{Velocity slotting (task 3.6)}

\texttt{scripts/compare\_slotting.py} draws random orders of 1 to 3 different items, each item as likely as its share
of the last 14 days' sales, and finds the bin each item would be picked from, as Inventory sends a picker: first expiry
first, then nearest. A walk costs twice the furthest \texttt{walk\_seq} it visits (Member 2's model, task 2.6). The
script then applies the plan and measures the same orders again (Table~\ref{tab:m3-slotting}).

\begin{table}[ht]
  \centering
  \footnotesize
  \begin{tabular}{@{}llrrrr@{}}
    \toprule
    Run & Orders measured & Orders & Before & After & Saving \\
    \midrule
    200 orders (seed 7)     & All orders         & 200     & 24.6 & 24.3 & 1\%  \\
                            & Ambient items only & 49      & 11.2 &  9.6 & 15\% \\
    \addlinespace[3pt]
    5{,}000 orders (seed 7) & All orders         & 5{,}000 & 24.1 & 23.9 & 1\%  \\
                            & Ambient items only & 1{,}317 & 10.6 &  9.9 & 6\%  \\
    \bottomrule
  \end{tabular}
  \caption{Average walk per order in bin steps, store 1, before and after the re-slotting plan was applied (13 batches
    moved, none skipped). The guide asks for 200 orders; 5,000 are less noisy.}
  \label{tab:m3-slotting}
\end{table}

The plan shortens the average walk by only 1\%. Orders of ambient items alone save 6\%; the 15\% of the 200-order run
is luck, and order seeds 1 to 10 gave between $-1$\% and 2\% for all orders and between $-9$\% and 15\% for ambient
ones. The plan for store 1 shows why:
\begin{itemize}
  \item \textbf{The fridges decide most walks.} Slotting re-sorts bins only within a temperature zone, and the chilled
    and frozen bins are at the far end of the walk (\texttt{walk\_seq} 13 to 18). Milk, curd and paneer sell 22 to 24
    a day each, 47\% of all units sold, so about three orders in four (3{,}683 of 5{,}000) include a chilled or frozen
    item and walk to the fridges whatever the plan.
  \item \textbf{The fastest ambient items were already at the front.} The demo layout had eggs (24.9 a day) and bread
    (14.8) in the first two bins; the plan only swaps them.
  \item \textbf{Equal sellers only trade places.} The three chilled items sell at almost the same rate, so curd and
    milk only swap the first two chilled bins. The other ambient items sell 3 to 5 a day, so their order is mostly noise.
\end{itemize}
In this layout the walk is decided by where the fridges stand, not by which shelf holds which item. A bigger saving
needs the fridges nearer the packing desk, which is a layout decision. Slotting should matter more in a real store,
with thousands of SKUs of very different speeds in each zone, than in a demo store with 16.

\section{Trade-off: for sale at putaway, not when the truck is done (HLD 13)}
\label{sec:m3-tradeoff}

The prototype first received a truck in one call: one transaction, all or nothing, and no record of the receipt. That
is the simplest correct design, because if anything fails, nothing happened. HLD 6.3 asks for more, each line for sale
the moment it is in its bin, and task 3.5 built that as five steps with a saved receipt. Building it showed what a
faster shelf costs.

\textbf{What it gains.} A fast seller that has run out is for sale as soon as its case is on the shelf, not after the
last case of the truck is scanned. The receipt also records what was ordered, scanned, damaged, refused, put away and
sent back, which the single call never kept.

\textbf{What it costs.}
\begin{itemize}
  \item \textbf{A third place for stock.} Units on the dock are neither on the shelf nor gone, and every rule had to say
    what happens to them: they still count as in transit, a receipt cannot close while they wait unless they are sent
    back, and a receipt can be cancelled only before anything from it is put away.
  \item \textbf{Races and retries the single call never had.} Two receivers could put away the same line, one could
    close the receipt while another scans, and a handheld could send a scan twice. Each step locks the receipt's row,
    units leave the dock through one conditional \texttt{UPDATE}, a bin's room is checked again after writing, and
    every write takes an \texttt{Idempotency-Key}, so a repeated scan counts once.
  \item \textbf{Edge cases keep appearing.} Review found that a line whose batch number is already in the store with
    another expiry date can be scanned onto the dock but never put away, so its receipt closes only by sending the line
    back. The fix, refusing such a scan at the dock, is waiting for review.
\end{itemize}

We kept the single call as well: \texttt{POST /v1/goods-receipts} runs the five steps in one transaction, so the ops
screen's \emph{Receive truck} button is still all or nothing, and both paths share one set of rules. We accepted a receipt with more states, and more
code to keep them consistent, in exchange for a shelf that refills line by line and a record of every truck.

\section{Limits and open items}

Table~\ref{tab:m3-limits} lists what the supply work leaves open, and what would close each item.

\begin{table}[!htb]
  \centering
  \footnotesize
  \begin{tabular}{@{}>{\raggedright\arraybackslash}p{0.27\linewidth}>{\raggedright\arraybackslash}p{0.36\linewidth}
                  >{\raggedright\arraybackslash}p{0.31\linewidth}@{}}
    \toprule
    Limit & What it means & Next step \\
    \midrule
    Daily forecasts, not hourly & The evening peak is not planned for & Spread each daily forecast over the hours by
      their share of sales \\
    \addlinespace[3pt]
    No festivals or weather & A festival day is missed, then repeated in the same weekday's forecast for four weeks
      & A festival calendar as an input \\
    \addlinespace[3pt]
    Reorder point from today's forecast only & Run on a busy Sunday, milk is sized for a Sunday & Size it from the
      forecasts of the lead-time days \\
    \addlinespace[3pt]
    The same cover for fresh and long-life food & Three days of cover above the reorder point, also for milk that keeps
      three days & Cover days that depend on shelf life \\
    \addlinespace[3pt]
    A batch number already in the store with another expiry date & It is scanned onto the dock but cannot be put away;
      only sending it back closes the receipt & Refuse it at the scan (pull request \#77, pending) \\
    \addlinespace[3pt]
    Slotting only within a zone; apply records the moves at once & 1\% shorter walks in this layout, and staff must
      carry the stock straight away & Fridges nearer the packing desk; items bought together side by side \\
    \bottomrule
  \end{tabular}
  \caption{Limits of the supply work and what would close them.}
  \label{tab:m3-limits}
\end{table}
. Needs: array, booktabs, tikz (arrows.meta, positioning,
% calc).
% member3-standalone.tex compiles these sections on their own.
%
% Numbers measured on main (8c83ed9) on Sunday 22 November 2026, on throwaway SQLite databases. The generated sales
% are drawn back from the day the script runs, so another day gives slightly different numbers. To re-run them, from
% the repository root with the project's virtual environment:
%   export DARKSTORE_DB=sqlite:///$(mktemp -d)/m3.db BACKGROUND_JOBS=off
%   python -m app.seed && python scripts/generate_sales.py            # demo data and 60 days of sales (seed 42)
%   python -c "from app.db import SessionLocal; from app.modules.forecasting.service import compare_models as c; db = SessionLocal(); [print(s, d, c(db, s, d)['wape']) for s in (1, 2) for d in (14, 28)]"
%     -> Table tab:m3-forecast. For the ranges in the text, repeat the three lines with generate_sales.py --seed 1 to 10.
%   python -c "from app.db import SessionLocal; from app.modules.forecasting import service as f; from app.modules.replenishment import service as r; db = SessionLocal(); print([x for x in f.run(db, 1) if x['sku_id'] == 1]); print(r.order_view(r.run(db, 1)['transfer_order'])); print(r.run(db, 1))"
%     -> the milk example in section sec:m3-replenishment (after the comparison: it changes the reorder points).
%   python scripts/compare_slotting.py --orders 200 --seed 7            # uses its own temporary database
%   python scripts/compare_slotting.py --orders 5000 --seed 7           # -> Table tab:m3-slotting; --seed 1 to 10 for the range

\section{Supply services (HLD 5, services 7, 8, 9, 6 and 2)}
\label{sec:m3-services}

Members 1 and 2 make sure that an order for something in stock goes well. Supply decides whether it is in stock in
the first place: the shelf should be full before the customer asks. Four services form a loop
(Figure~\ref{fig:m3-loop}), and Assortment sits above it, deciding which SKUs a store carries at all. None of the five
writes a stock quantity: Inbound asks Inventory to receive and Slotting asks it to move, so Member 1's three stock
records always agree. Forecasting writes only two planning settings on each stock row, the reorder point and the
target stock.

\begin{figure}[ht]
  \centering
  \begin{tikzpicture}[
      >={Stealth[length=2.2mm]}, node distance=17mm,
      svc/.style={draw=black!45, rounded corners=3pt, minimum width=25mm, minimum height=11mm, align=center,
                  font=\bfseries\small},
      shelf/.style={svc, draw=blue!70!black, fill=blue!8, line width=0.9pt, minimum width=62mm},
      edgelabel/.style={font=\footnotesize, text=black!65, align=center}]
    \node[svc] (fc) {Forecasting\\[-1pt]{\mdseries\scriptsize service 7}};
    \node[svc, right=of fc] (rp) {Replenishment\\[-1pt]{\mdseries\scriptsize service 8}};
    \node[svc, right=of rp] (ib) {Inbound\\[-1pt]{\mdseries\scriptsize service 9}};
    \node[svc, right=of ib] (sl) {Slotting\\[-1pt]{\mdseries\scriptsize service 6}};
    \node[svc, above=10mm of rp] (as) {Assortment\\[-1pt]{\mdseries\scriptsize service 2}};
    \node[shelf, below=16mm of $(rp)!0.5!(ib)$] (shelf) {Stock in bins\\[-1pt]{\mdseries\scriptsize Inventory
      (Member 1) writes every change}};

    \draw[->, black!55] (fc) -- node[edgelabel, above] {reorder\\point} (rp);
    \draw[->, black!55] (rp) -- node[edgelabel, above] {transfer\\order} (ib);
    \draw[->, black!55] (sl) -- node[edgelabel, above] {bin per\\line} (ib);
    \draw[->, black!55] (as) -- node[edgelabel, right] {SKUs carried} (rp);
    \draw[->, black!55] (ib.south) -- node[edgelabel, right] {put away} (ib.south |- shelf.north);
    \draw[->, black!55] (sl.south) |- node[edgelabel, pos=0.3, right] {moves} (shelf.east);
    \draw[->, black!55] (shelf.west) -| node[edgelabel, pos=0.75, left] {sales\\(ledger)} (fc.south);
  \end{tikzpicture}
  \caption{The supply loop. Only Inventory changes stock: Inbound and Slotting ask it to receive and to move units.}
  \label{fig:m3-loop}
\end{figure}

\subsection{Forecasting (service 7)}

Forecasting predicts each SKU's daily sales and turns them into a reorder point (FR-7). Sales come from the
\texttt{sale} rows of the ledger, in whole days ending yesterday, because today is only partly over.

\begin{itemize}
  \item \textbf{Three models with one shape.} Each forecasts a day from the days before it: the \emph{moving average}
    of the last 7 days, the \emph{weekday average} of the same weekday in the last 4 weeks, and \emph{exponential
    smoothing} ($\alpha = 0.3$), in which each earlier day counts 0.7 times as much as the day after it. Days before a
    SKU's first sale are unknown, not zero, so a new SKU is not forecast as if it had sold nothing for weeks.
  \item \textbf{A measuring stick first.} The backtest replays the last 14 days, forecasting each day from the days
    before it, and reports the error over every SKU and day,
    \[ \mathrm{WAPE} = \frac{\sum |\textrm{actual} - \textrm{forecast}|}{\sum \textrm{actual}}, \]
    for which HLD 12 asks 30\% or lower. Weighting by volume stops slow sellers swamping the score. A day counts only
    when the SKU has a week of sales before it, so on the demo data's single week of sales the answer is ``not
    available'', not a misleading number. A comparison runs all three models on the same days, and the setting
    \texttt{FORECAST\_MODEL} picks the one in use: the weekday average, which won (section~\ref{sec:m3-results}).
  \item \textbf{Safety stock from the model's errors.} The reorder point uses the spread of the model's \emph{errors}
    over the last 14 days, not the spread of raw sales: safety stock should cover what the forecast gets wrong, not
    variation it already predicts, such as busier weekends.
\end{itemize}
A manager runs it for a store (\texttt{POST /v1/forecasting/run}), and each forecast is saved with its model's name.

\subsection{Replenishment (service 8)}

Replenishment orders from the mother hub when stock runs low (FR-8), in whole cases and within each zone's free bin
space. It runs every 15 minutes for every open store, and whenever a manager asks; section~\ref{sec:m3-replenishment}
follows one run. It also publishes \texttt{stock.low} when the stock not held for customers falls below the reorder point, lets
a manager stop reordering one SKU so that it runs down before it is dropped, and reports the five supply KPIs of HLD 12
(stock-outs, days of cover, wastage, deliveries on time and forecast WAPE) against their targets for Member 4's ops
dashboard. The hub itself is a stand-in: a button plays its warehouse system saying ``the truck has left''.

\subsection{Inbound (service 9)}

Inbound receives a truck in the five steps of HLD 6.3 (section~\ref{sec:m3-inbound}). Each truck gets a saved goods
receipt whose lines record, per SKU and batch, the units scanned, damaged, refused and put away. Putting a line away
calls \texttt{inventory.receive} (Member 1), which writes the bin, the totals row and the ledger in one transaction.

\subsection{Slotting (service 6)}

Slotting decides which bin a SKU lives in. Each bin has a zone (ambient, chilled or frozen), a capacity and a
\texttt{walk\_seq}, its place on the picker's route, where bin 1 is next to the packing desk. HLD service 6 ``places
fast movers near the packing station'':
\begin{itemize}
  \item \textbf{The plan} (look only) ranks each zone's SKUs by units sold a day over the last 14 days and gives the
    fastest the bin nearest the desk, the next one the next bin, and so on. A SKU that sold nothing stays where it is
    unless a seller needs its bin, since moving it would not shorten any walk.
  \item \textbf{Apply} (managers) locks the store's bins and moves each batch through \texttt{inventory.move}: two
    ledger rows, nothing gained or lost. A batch moves only when its new bin has room at that moment, so no bin is ever
    over capacity, and a batch that a picker has been sent to stays where it is.
  \item \textbf{New deliveries} go to the SKU's planned bin first, then to bins already holding it, then to the
    emptiest bins of its zone, never past capacity, so the shelf re-sorts itself as stock arrives.
\end{itemize}
A SKU never leaves its own temperature zone; section~\ref{sec:m3-results} shows what that costs in walking.

\subsection{Assortment (service 2)}

Assortment decides which SKUs each store carries. Adding a SKU creates its stock row with zero stock, so
Replenishment starts ordering it; dropping one is refused while any unit is on the shelf, held for a customer or on its
way. From the last 30 days of sales it suggests SKUs to \emph{add} (selling in other stores, highest demand first),
SKUs to \emph{drop} (stock on hand 30 days ago and not one unit sold since) and the \emph{space} left in each zone
(bins against SKUs carried, so nobody adds a fifth chilled SKU to four chilled bins). A manager decides.

\section{Replenishment (HLD 6.2)}
\label{sec:m3-replenishment}

HLD 6.2 keeps every SKU above a reorder point. The HLD counts in hours and the prototype in days, with a lead time $L$
of 1 day, $z = 1.65$ (stock should run out in at most 5\% of cycles) and 3 days of cover:
\[
  \textrm{reorder point} = \bigl\lceil dL + z\sigma\sqrt{L}\, \bigr\rceil \qquad\qquad
  \textrm{target stock} = \textrm{reorder point} + \lceil 3d \rceil
\]
Here $d$ is the forecast daily demand and $\sigma$ the spread of the forecast's errors. The first part of the reorder
point covers the sales expected while the truck is on its way; the second is safety stock for a busier day.

\begin{enumerate}
  \item \textbf{Check.} Each SKU's \emph{position}, on hand $-$ held $+$ in transit, is compared with its reorder
    point. Counting what is already on its way is what makes the 15-minute timer harmless: without it, every run would
    order the same items again until the truck came.
  \item \textbf{Size.} A SKU below its reorder point needs target stock $-$ position, rounded up to whole cases
    (12 milk, 6 eggs, 10 bread).
  \item \textbf{Fit.} A zone's free space is its bins' capacity minus the units stored and the units on their way.
    When the lines do not fit, the most urgent go first: lowest days of cover (what customers can buy $\div$ daily
    demand), then the faster seller. The last line is cut to the cases that fit, and every cut is reported, also when
    the timer made it.
  \item \textbf{Order.} One transfer order takes every line (created, dispatched, received). Only one run happens at a
    time in an app process, so the timer and a manager's button cannot both order the same stock.
\end{enumerate}

\paragraph{One run, with real numbers.} On the generated sales of section~\ref{sec:m3-results}, run on Sunday
22 November, milk in store 1 is forecast at 32.25 a day: the mean of the last four Sundays, one of them a festival day.
The spread of its errors is 6.5, so the reorder point is $\lceil 32.25 + 1.65 \times 6.5 \rceil = 43$ and the target
$43 + \lceil 3 \times 32.25 \rceil = 140$. With 30 on hand, none held and none on the way, the run orders
$140 - 30 = 110$, rounded up to 10 cases of 12: 120. The same transfer order took bread (4 left: 9 cases of 10), eggs,
curd, paneer and ice cream, and nothing was cut: the three chilled lines needed 362 of the 510 free chilled units. A
second run straight after ordered nothing.

The example also shows a weakness: the reorder point is sized from \emph{today's} forecast only. Run on a Sunday after
a festival Sunday, milk is sized for 32 a day; run on a weekday, it would be sized for a weekday.

\section{Inbound (HLD 6.3)}
\label{sec:m3-inbound}

HLD 6.3 receives a truck in five steps and asks that ``each item becomes sellable as soon as its bin scan is
confirmed, not when the whole truck is done''. In the prototype only step 4 changes stock.

\begin{enumerate}
  \item \textbf{Open.} A receipt is saved for the transfer order. The order is locked first, so a double click on
    \emph{Receive truck} opens one receipt, not two.
  \item \textbf{Scan.} Each case is scanned with its batch number and expiry date, which first-expiry-first-out
    picking relies on. Good units wait on the dock, not yet for sale. Damaged units, and a batch with less than half
    its shelf life left (milk, which keeps 3 days, with 1 day to go), are recorded but never stocked.
  \item \textbf{Putaway tasks.} Slotting names a bin, or several, for each line on the dock, never past capacity.
  \item \textbf{Put away.} The units go into the bin through \texttt{inventory.receive} and are \emph{for sale as soon
    as this commits}, whatever is still on the dock. They stop counting as in transit at the same moment.
  \item \textbf{Close.} The transfer order is marked received, and every SKU is reported as ordered against received,
    with damaged, refused and sent-back units. Closing is refused while good units wait on the dock, unless the
    receiver sends them back to the hub.
\end{enumerate}

A receipt opened by mistake can be cancelled until anything from it is put away, and the real truck can still be
received. The tests check the property HLD 6.3 asks for: after a scan, what customers can buy is unchanged; after one
line is put away, that SKU rises and no other does.

\section{Measurements}
\label{sec:m3-results}

Both measurements run the project's own code on a throwaway copy of the demo data with 60 days of generated sales
(\texttt{scripts/generate\_sales.py}, task 3.1, seed 42): for every SKU and day, a random (Poisson) number of units
around a base demand per category (20 a day for dairy, including eggs, 12 for bread, 4 for staples, 2 for frozen
food, 3 for the rest), 35\% more at weekends and twice as much on two festival days. They were measured on main on
22 November 2026; the commands are at the top of \texttt{member3-supply.tex}. The history is drawn back from the day
it is generated, so another day gives slightly different numbers.

\subsection{Forecast models (task 3.3)}

\begin{table}[ht]
  \centering
  \footnotesize
  \begin{tabular}{@{}lrrrr@{}}
    \toprule
    Model & \shortstack[r]{Store 1\\14 days} & \shortstack[r]{Store 2\\14 days} & \shortstack[r]{Store 1\\28 days}
          & \shortstack[r]{Store 2\\28 days} \\
    \midrule
    Moving average, last 7 days (the baseline) & 34.4\% & 36.1\% & 30.7\% & 32.4\% \\
    Weekday average, last 4 same weekdays      & \textbf{29.3\%} & \textbf{29.2\%} & \textbf{28.9\%} & \textbf{27.8\%} \\
    Exponential smoothing, $\alpha = 0.3$      & 35.9\% & 39.2\% & 32.4\% & 34.3\% \\
    \bottomrule
  \end{tabular}
  \caption{Forecast error (WAPE, lower is better) of each model over the last 14 and 28 days, every SKU of the store
    (\texttt{GET /v1/forecasting/compare}). HLD 12 target: 30\% or lower.}
  \label{tab:m3-forecast}
\end{table}

The weekday average wins in both stores and both windows (Table~\ref{tab:m3-forecast}). Over 14 days it is 5 to 7
points better than the baseline and the only model within the 30\% target, so it is the model in use. With ten other
random seeds for the history (40 more store-and-window cases) it was again the best every time, at 26.0 to 32.3\%
against 30.5 to 39.1\% for the moving average and 32.1 to 41.1\% for smoothing, and it met the target in 28 of the 40.

\begin{itemize}
  \item \textbf{Why it wins.} A 7-day moving average forecasts every day as the week's average, so with weekends 35\%
    busier it is about 10\% too high on weekdays and 19\% too low at weekends. The weekday average forecasts a
    Saturday from Saturdays. Smoothing does worst, because it leans on the last two or three days: a Monday is
    forecast mostly from the weekend, and a Saturday from the weekdays.
  \item \textbf{What is left is mostly noise.} For a SKU that sells 3 a day, one unit off is already a third of the
    day's sales, so the slow sellers keep the store near 30\% whatever the model. The 14-day window also starts on a
    festival day, which every model misses; the weekday average then carries it into the same weekday's forecast for
    four weeks, as in the milk example of section~\ref{sec:m3-replenishment}.
\end{itemize}

\subsection{Velocity slotting (task 3.6)}

\texttt{scripts/compare\_slotting.py} draws random orders of 1 to 3 different items, each item as likely as its share
of the last 14 days' sales, and finds the bin each item would be picked from, as Inventory sends a picker: first expiry
first, then nearest. A walk costs twice the furthest \texttt{walk\_seq} it visits (Member 2's model, task 2.6). The
script then applies the plan and measures the same orders again (Table~\ref{tab:m3-slotting}).

\begin{table}[ht]
  \centering
  \footnotesize
  \begin{tabular}{@{}llrrrr@{}}
    \toprule
    Run & Orders measured & Orders & Before & After & Saving \\
    \midrule
    200 orders (seed 7)     & All orders         & 200     & 24.6 & 24.3 & 1\%  \\
                            & Ambient items only & 49      & 11.2 &  9.6 & 15\% \\
    \addlinespace[3pt]
    5{,}000 orders (seed 7) & All orders         & 5{,}000 & 24.1 & 23.9 & 1\%  \\
                            & Ambient items only & 1{,}317 & 10.6 &  9.9 & 6\%  \\
    \bottomrule
  \end{tabular}
  \caption{Average walk per order in bin steps, store 1, before and after the re-slotting plan was applied (13 batches
    moved, none skipped). The guide asks for 200 orders; 5,000 are less noisy.}
  \label{tab:m3-slotting}
\end{table}

The plan shortens the average walk by only 1\%. Orders of ambient items alone save 6\%; the 15\% of the 200-order run
is luck, and order seeds 1 to 10 gave between $-1$\% and 2\% for all orders and between $-9$\% and 15\% for ambient
ones. The plan for store 1 shows why:
\begin{itemize}
  \item \textbf{The fridges decide most walks.} Slotting re-sorts bins only within a temperature zone, and the chilled
    and frozen bins are at the far end of the walk (\texttt{walk\_seq} 13 to 18). Milk, curd and paneer sell 22 to 24
    a day each, 47\% of all units sold, so about three orders in four (3{,}683 of 5{,}000) include a chilled or frozen
    item and walk to the fridges whatever the plan.
  \item \textbf{The fastest ambient items were already at the front.} The demo layout had eggs (24.9 a day) and bread
    (14.8) in the first two bins; the plan only swaps them.
  \item \textbf{Equal sellers only trade places.} The three chilled items sell at almost the same rate, so curd and
    milk only swap the first two chilled bins. The other ambient items sell 3 to 5 a day, so their order is mostly noise.
\end{itemize}
In this layout the walk is decided by where the fridges stand, not by which shelf holds which item. A bigger saving
needs the fridges nearer the packing desk, which is a layout decision. Slotting should matter more in a real store,
with thousands of SKUs of very different speeds in each zone, than in a demo store with 16.

\section{Trade-off: for sale at putaway, not when the truck is done (HLD 13)}
\label{sec:m3-tradeoff}

The prototype first received a truck in one call: one transaction, all or nothing, and no record of the receipt. That
is the simplest correct design, because if anything fails, nothing happened. HLD 6.3 asks for more, each line for sale
the moment it is in its bin, and task 3.5 built that as five steps with a saved receipt. Building it showed what a
faster shelf costs.

\textbf{What it gains.} A fast seller that has run out is for sale as soon as its case is on the shelf, not after the
last case of the truck is scanned. The receipt also records what was ordered, scanned, damaged, refused, put away and
sent back, which the single call never kept.

\textbf{What it costs.}
\begin{itemize}
  \item \textbf{A third place for stock.} Units on the dock are neither on the shelf nor gone, and every rule had to say
    what happens to them: they still count as in transit, a receipt cannot close while they wait unless they are sent
    back, and a receipt can be cancelled only before anything from it is put away.
  \item \textbf{Races and retries the single call never had.} Two receivers could put away the same line, one could
    close the receipt while another scans, and a handheld could send a scan twice. Each step locks the receipt's row,
    units leave the dock through one conditional \texttt{UPDATE}, a bin's room is checked again after writing, and
    every write takes an \texttt{Idempotency-Key}, so a repeated scan counts once.
  \item \textbf{Edge cases keep appearing.} Review found that a line whose batch number is already in the store with
    another expiry date can be scanned onto the dock but never put away, so its receipt closes only by sending the line
    back. The fix, refusing such a scan at the dock, is waiting for review.
\end{itemize}

We kept the single call as well: \texttt{POST /v1/goods-receipts} runs the five steps in one transaction, so the ops
screen's \emph{Receive truck} button is still all or nothing, and both paths share one set of rules. We accepted a receipt with more states, and more
code to keep them consistent, in exchange for a shelf that refills line by line and a record of every truck.

\section{Limits and open items}

Table~\ref{tab:m3-limits} lists what the supply work leaves open, and what would close each item.

\begin{table}[!htb]
  \centering
  \footnotesize
  \begin{tabular}{@{}>{\raggedright\arraybackslash}p{0.27\linewidth}>{\raggedright\arraybackslash}p{0.36\linewidth}
                  >{\raggedright\arraybackslash}p{0.31\linewidth}@{}}
    \toprule
    Limit & What it means & Next step \\
    \midrule
    Daily forecasts, not hourly & The evening peak is not planned for & Spread each daily forecast over the hours by
      their share of sales \\
    \addlinespace[3pt]
    No festivals or weather & A festival day is missed, then repeated in the same weekday's forecast for four weeks
      & A festival calendar as an input \\
    \addlinespace[3pt]
    Reorder point from today's forecast only & Run on a busy Sunday, milk is sized for a Sunday & Size it from the
      forecasts of the lead-time days \\
    \addlinespace[3pt]
    The same cover for fresh and long-life food & Three days of cover above the reorder point, also for milk that keeps
      three days & Cover days that depend on shelf life \\
    \addlinespace[3pt]
    A batch number already in the store with another expiry date & It is scanned onto the dock but cannot be put away;
      only sending it back closes the receipt & Refuse it at the scan (pull request \#77, pending) \\
    \addlinespace[3pt]
    Slotting only within a zone; apply records the moves at once & 1\% shorter walks in this layout, and staff must
      carry the stock straight away & Fridges nearer the packing desk; items bought together side by side \\
    \bottomrule
  \end{tabular}
  \caption{Limits of the supply work and what would close them.}
  \label{tab:m3-limits}
\end{table}
. Needs: array, booktabs, tikz (arrows.meta, positioning,
% calc).
% member3-standalone.tex compiles these sections on their own.
%
% Numbers measured on main (8c83ed9) on Sunday 22 November 2026, on throwaway SQLite databases. The generated sales
% are drawn back from the day the script runs, so another day gives slightly different numbers. To re-run them, from
% the repository root with the project's virtual environment:
%   export DARKSTORE_DB=sqlite:///$(mktemp -d)/m3.db BACKGROUND_JOBS=off
%   python -m app.seed && python scripts/generate_sales.py            # demo data and 60 days of sales (seed 42)
%   python -c "from app.db import SessionLocal; from app.modules.forecasting.service import compare_models as c; db = SessionLocal(); [print(s, d, c(db, s, d)['wape']) for s in (1, 2) for d in (14, 28)]"
%     -> Table tab:m3-forecast. For the ranges in the text, repeat the three lines with generate_sales.py --seed 1 to 10.
%   python -c "from app.db import SessionLocal; from app.modules.forecasting import service as f; from app.modules.replenishment import service as r; db = SessionLocal(); print([x for x in f.run(db, 1) if x['sku_id'] == 1]); print(r.order_view(r.run(db, 1)['transfer_order'])); print(r.run(db, 1))"
%     -> the milk example in section sec:m3-replenishment (after the comparison: it changes the reorder points).
%   python scripts/compare_slotting.py --orders 200 --seed 7            # uses its own temporary database
%   python scripts/compare_slotting.py --orders 5000 --seed 7           # -> Table tab:m3-slotting; --seed 1 to 10 for the range

\section{Supply services (HLD 5, services 7, 8, 9, 6 and 2)}
\label{sec:m3-services}

Members 1 and 2 make sure that an order for something in stock goes well. Supply decides whether it is in stock in
the first place: the shelf should be full before the customer asks. Four services form a loop
(Figure~\ref{fig:m3-loop}), and Assortment sits above it, deciding which SKUs a store carries at all. None of the five
writes a stock quantity: Inbound asks Inventory to receive and Slotting asks it to move, so Member 1's three stock
records always agree. Forecasting writes only two planning settings on each stock row, the reorder point and the
target stock.

\begin{figure}[ht]
  \centering
  \begin{tikzpicture}[
      >={Stealth[length=2.2mm]}, node distance=17mm,
      svc/.style={draw=black!45, rounded corners=3pt, minimum width=25mm, minimum height=11mm, align=center,
                  font=\bfseries\small},
      shelf/.style={svc, draw=blue!70!black, fill=blue!8, line width=0.9pt, minimum width=62mm},
      edgelabel/.style={font=\footnotesize, text=black!65, align=center}]
    \node[svc] (fc) {Forecasting\\[-1pt]{\mdseries\scriptsize service 7}};
    \node[svc, right=of fc] (rp) {Replenishment\\[-1pt]{\mdseries\scriptsize service 8}};
    \node[svc, right=of rp] (ib) {Inbound\\[-1pt]{\mdseries\scriptsize service 9}};
    \node[svc, right=of ib] (sl) {Slotting\\[-1pt]{\mdseries\scriptsize service 6}};
    \node[svc, above=10mm of rp] (as) {Assortment\\[-1pt]{\mdseries\scriptsize service 2}};
    \node[shelf, below=16mm of $(rp)!0.5!(ib)$] (shelf) {Stock in bins\\[-1pt]{\mdseries\scriptsize Inventory
      (Member 1) writes every change}};

    \draw[->, black!55] (fc) -- node[edgelabel, above] {reorder\\point} (rp);
    \draw[->, black!55] (rp) -- node[edgelabel, above] {transfer\\order} (ib);
    \draw[->, black!55] (sl) -- node[edgelabel, above] {bin per\\line} (ib);
    \draw[->, black!55] (as) -- node[edgelabel, right] {SKUs carried} (rp);
    \draw[->, black!55] (ib.south) -- node[edgelabel, right] {put away} (ib.south |- shelf.north);
    \draw[->, black!55] (sl.south) |- node[edgelabel, pos=0.3, right] {moves} (shelf.east);
    \draw[->, black!55] (shelf.west) -| node[edgelabel, pos=0.75, left] {sales\\(ledger)} (fc.south);
  \end{tikzpicture}
  \caption{The supply loop. Only Inventory changes stock: Inbound and Slotting ask it to receive and to move units.}
  \label{fig:m3-loop}
\end{figure}

\subsection{Forecasting (service 7)}

Forecasting predicts each SKU's daily sales and turns them into a reorder point (FR-7). Sales come from the
\texttt{sale} rows of the ledger, in whole days ending yesterday, because today is only partly over.

\begin{itemize}
  \item \textbf{Three models with one shape.} Each forecasts a day from the days before it: the \emph{moving average}
    of the last 7 days, the \emph{weekday average} of the same weekday in the last 4 weeks, and \emph{exponential
    smoothing} ($\alpha = 0.3$), in which each earlier day counts 0.7 times as much as the day after it. Days before a
    SKU's first sale are unknown, not zero, so a new SKU is not forecast as if it had sold nothing for weeks.
  \item \textbf{A measuring stick first.} The backtest replays the last 14 days, forecasting each day from the days
    before it, and reports the error over every SKU and day,
    \[ \mathrm{WAPE} = \frac{\sum |\textrm{actual} - \textrm{forecast}|}{\sum \textrm{actual}}, \]
    for which HLD 12 asks 30\% or lower. Weighting by volume stops slow sellers swamping the score. A day counts only
    when the SKU has a week of sales before it, so on the demo data's single week of sales the answer is ``not
    available'', not a misleading number. A comparison runs all three models on the same days, and the setting
    \texttt{FORECAST\_MODEL} picks the one in use: the weekday average, which won (section~\ref{sec:m3-results}).
  \item \textbf{Safety stock from the model's errors.} The reorder point uses the spread of the model's \emph{errors}
    over the last 14 days, not the spread of raw sales: safety stock should cover what the forecast gets wrong, not
    variation it already predicts, such as busier weekends.
\end{itemize}
A manager runs it for a store (\texttt{POST /v1/forecasting/run}), and each forecast is saved with its model's name.

\subsection{Replenishment (service 8)}

Replenishment orders from the mother hub when stock runs low (FR-8), in whole cases and within each zone's free bin
space. It runs every 15 minutes for every open store, and whenever a manager asks; section~\ref{sec:m3-replenishment}
follows one run. It also publishes \texttt{stock.low} when the stock not held for customers falls below the reorder point, lets
a manager stop reordering one SKU so that it runs down before it is dropped, and reports the five supply KPIs of HLD 12
(stock-outs, days of cover, wastage, deliveries on time and forecast WAPE) against their targets for Member 4's ops
dashboard. The hub itself is a stand-in: a button plays its warehouse system saying ``the truck has left''.

\subsection{Inbound (service 9)}

Inbound receives a truck in the five steps of HLD 6.3 (section~\ref{sec:m3-inbound}). Each truck gets a saved goods
receipt whose lines record, per SKU and batch, the units scanned, damaged, refused and put away. Putting a line away
calls \texttt{inventory.receive} (Member 1), which writes the bin, the totals row and the ledger in one transaction.

\subsection{Slotting (service 6)}

Slotting decides which bin a SKU lives in. Each bin has a zone (ambient, chilled or frozen), a capacity and a
\texttt{walk\_seq}, its place on the picker's route, where bin 1 is next to the packing desk. HLD service 6 ``places
fast movers near the packing station'':
\begin{itemize}
  \item \textbf{The plan} (look only) ranks each zone's SKUs by units sold a day over the last 14 days and gives the
    fastest the bin nearest the desk, the next one the next bin, and so on. A SKU that sold nothing stays where it is
    unless a seller needs its bin, since moving it would not shorten any walk.
  \item \textbf{Apply} (managers) locks the store's bins and moves each batch through \texttt{inventory.move}: two
    ledger rows, nothing gained or lost. A batch moves only when its new bin has room at that moment, so no bin is ever
    over capacity, and a batch that a picker has been sent to stays where it is.
  \item \textbf{New deliveries} go to the SKU's planned bin first, then to bins already holding it, then to the
    emptiest bins of its zone, never past capacity, so the shelf re-sorts itself as stock arrives.
\end{itemize}
A SKU never leaves its own temperature zone; section~\ref{sec:m3-results} shows what that costs in walking.

\subsection{Assortment (service 2)}

Assortment decides which SKUs each store carries. Adding a SKU creates its stock row with zero stock, so
Replenishment starts ordering it; dropping one is refused while any unit is on the shelf, held for a customer or on its
way. From the last 30 days of sales it suggests SKUs to \emph{add} (selling in other stores, highest demand first),
SKUs to \emph{drop} (stock on hand 30 days ago and not one unit sold since) and the \emph{space} left in each zone
(bins against SKUs carried, so nobody adds a fifth chilled SKU to four chilled bins). A manager decides.

\section{Replenishment (HLD 6.2)}
\label{sec:m3-replenishment}

HLD 6.2 keeps every SKU above a reorder point. The HLD counts in hours and the prototype in days, with a lead time $L$
of 1 day, $z = 1.65$ (stock should run out in at most 5\% of cycles) and 3 days of cover:
\[
  \textrm{reorder point} = \bigl\lceil dL + z\sigma\sqrt{L}\, \bigr\rceil \qquad\qquad
  \textrm{target stock} = \textrm{reorder point} + \lceil 3d \rceil
\]
Here $d$ is the forecast daily demand and $\sigma$ the spread of the forecast's errors. The first part of the reorder
point covers the sales expected while the truck is on its way; the second is safety stock for a busier day.

\begin{enumerate}
  \item \textbf{Check.} Each SKU's \emph{position}, on hand $-$ held $+$ in transit, is compared with its reorder
    point. Counting what is already on its way is what makes the 15-minute timer harmless: without it, every run would
    order the same items again until the truck came.
  \item \textbf{Size.} A SKU below its reorder point needs target stock $-$ position, rounded up to whole cases
    (12 milk, 6 eggs, 10 bread).
  \item \textbf{Fit.} A zone's free space is its bins' capacity minus the units stored and the units on their way.
    When the lines do not fit, the most urgent go first: lowest days of cover (what customers can buy $\div$ daily
    demand), then the faster seller. The last line is cut to the cases that fit, and every cut is reported, also when
    the timer made it.
  \item \textbf{Order.} One transfer order takes every line (created, dispatched, received). Only one run happens at a
    time in an app process, so the timer and a manager's button cannot both order the same stock.
\end{enumerate}

\paragraph{One run, with real numbers.} On the generated sales of section~\ref{sec:m3-results}, run on Sunday
22 November, milk in store 1 is forecast at 32.25 a day: the mean of the last four Sundays, one of them a festival day.
The spread of its errors is 6.5, so the reorder point is $\lceil 32.25 + 1.65 \times 6.5 \rceil = 43$ and the target
$43 + \lceil 3 \times 32.25 \rceil = 140$. With 30 on hand, none held and none on the way, the run orders
$140 - 30 = 110$, rounded up to 10 cases of 12: 120. The same transfer order took bread (4 left: 9 cases of 10), eggs,
curd, paneer and ice cream, and nothing was cut: the three chilled lines needed 362 of the 510 free chilled units. A
second run straight after ordered nothing.

The example also shows a weakness: the reorder point is sized from \emph{today's} forecast only. Run on a Sunday after
a festival Sunday, milk is sized for 32 a day; run on a weekday, it would be sized for a weekday.

\section{Inbound (HLD 6.3)}
\label{sec:m3-inbound}

HLD 6.3 receives a truck in five steps and asks that ``each item becomes sellable as soon as its bin scan is
confirmed, not when the whole truck is done''. In the prototype only step 4 changes stock.

\begin{enumerate}
  \item \textbf{Open.} A receipt is saved for the transfer order. The order is locked first, so a double click on
    \emph{Receive truck} opens one receipt, not two.
  \item \textbf{Scan.} Each case is scanned with its batch number and expiry date, which first-expiry-first-out
    picking relies on. Good units wait on the dock, not yet for sale. Damaged units, and a batch with less than half
    its shelf life left (milk, which keeps 3 days, with 1 day to go), are recorded but never stocked.
  \item \textbf{Putaway tasks.} Slotting names a bin, or several, for each line on the dock, never past capacity.
  \item \textbf{Put away.} The units go into the bin through \texttt{inventory.receive} and are \emph{for sale as soon
    as this commits}, whatever is still on the dock. They stop counting as in transit at the same moment.
  \item \textbf{Close.} The transfer order is marked received, and every SKU is reported as ordered against received,
    with damaged, refused and sent-back units. Closing is refused while good units wait on the dock, unless the
    receiver sends them back to the hub.
\end{enumerate}

A receipt opened by mistake can be cancelled until anything from it is put away, and the real truck can still be
received. The tests check the property HLD 6.3 asks for: after a scan, what customers can buy is unchanged; after one
line is put away, that SKU rises and no other does.

\section{Measurements}
\label{sec:m3-results}

Both measurements run the project's own code on a throwaway copy of the demo data with 60 days of generated sales
(\texttt{scripts/generate\_sales.py}, task 3.1, seed 42): for every SKU and day, a random (Poisson) number of units
around a base demand per category (20 a day for dairy, including eggs, 12 for bread, 4 for staples, 2 for frozen
food, 3 for the rest), 35\% more at weekends and twice as much on two festival days. They were measured on main on
22 November 2026; the commands are at the top of \texttt{member3-supply.tex}. The history is drawn back from the day
it is generated, so another day gives slightly different numbers.

\subsection{Forecast models (task 3.3)}

\begin{table}[ht]
  \centering
  \footnotesize
  \begin{tabular}{@{}lrrrr@{}}
    \toprule
    Model & \shortstack[r]{Store 1\\14 days} & \shortstack[r]{Store 2\\14 days} & \shortstack[r]{Store 1\\28 days}
          & \shortstack[r]{Store 2\\28 days} \\
    \midrule
    Moving average, last 7 days (the baseline) & 34.4\% & 36.1\% & 30.7\% & 32.4\% \\
    Weekday average, last 4 same weekdays      & \textbf{29.3\%} & \textbf{29.2\%} & \textbf{28.9\%} & \textbf{27.8\%} \\
    Exponential smoothing, $\alpha = 0.3$      & 35.9\% & 39.2\% & 32.4\% & 34.3\% \\
    \bottomrule
  \end{tabular}
  \caption{Forecast error (WAPE, lower is better) of each model over the last 14 and 28 days, every SKU of the store
    (\texttt{GET /v1/forecasting/compare}). HLD 12 target: 30\% or lower.}
  \label{tab:m3-forecast}
\end{table}

The weekday average wins in both stores and both windows (Table~\ref{tab:m3-forecast}). Over 14 days it is 5 to 7
points better than the baseline and the only model within the 30\% target, so it is the model in use. With ten other
random seeds for the history (40 more store-and-window cases) it was again the best every time, at 26.0 to 32.3\%
against 30.5 to 39.1\% for the moving average and 32.1 to 41.1\% for smoothing, and it met the target in 28 of the 40.

\begin{itemize}
  \item \textbf{Why it wins.} A 7-day moving average forecasts every day as the week's average, so with weekends 35\%
    busier it is about 10\% too high on weekdays and 19\% too low at weekends. The weekday average forecasts a
    Saturday from Saturdays. Smoothing does worst, because it leans on the last two or three days: a Monday is
    forecast mostly from the weekend, and a Saturday from the weekdays.
  \item \textbf{What is left is mostly noise.} For a SKU that sells 3 a day, one unit off is already a third of the
    day's sales, so the slow sellers keep the store near 30\% whatever the model. The 14-day window also starts on a
    festival day, which every model misses; the weekday average then carries it into the same weekday's forecast for
    four weeks, as in the milk example of section~\ref{sec:m3-replenishment}.
\end{itemize}

\subsection{Velocity slotting (task 3.6)}

\texttt{scripts/compare\_slotting.py} draws random orders of 1 to 3 different items, each item as likely as its share
of the last 14 days' sales, and finds the bin each item would be picked from, as Inventory sends a picker: first expiry
first, then nearest. A walk costs twice the furthest \texttt{walk\_seq} it visits (Member 2's model, task 2.6). The
script then applies the plan and measures the same orders again (Table~\ref{tab:m3-slotting}).

\begin{table}[ht]
  \centering
  \footnotesize
  \begin{tabular}{@{}llrrrr@{}}
    \toprule
    Run & Orders measured & Orders & Before & After & Saving \\
    \midrule
    200 orders (seed 7)     & All orders         & 200     & 24.6 & 24.3 & 1\%  \\
                            & Ambient items only & 49      & 11.2 &  9.6 & 15\% \\
    \addlinespace[3pt]
    5{,}000 orders (seed 7) & All orders         & 5{,}000 & 24.1 & 23.9 & 1\%  \\
                            & Ambient items only & 1{,}317 & 10.6 &  9.9 & 6\%  \\
    \bottomrule
  \end{tabular}
  \caption{Average walk per order in bin steps, store 1, before and after the re-slotting plan was applied (13 batches
    moved, none skipped). The guide asks for 200 orders; 5,000 are less noisy.}
  \label{tab:m3-slotting}
\end{table}

The plan shortens the average walk by only 1\%. Orders of ambient items alone save 6\%; the 15\% of the 200-order run
is luck, and order seeds 1 to 10 gave between $-1$\% and 2\% for all orders and between $-9$\% and 15\% for ambient
ones. The plan for store 1 shows why:
\begin{itemize}
  \item \textbf{The fridges decide most walks.} Slotting re-sorts bins only within a temperature zone, and the chilled
    and frozen bins are at the far end of the walk (\texttt{walk\_seq} 13 to 18). Milk, curd and paneer sell 22 to 24
    a day each, 47\% of all units sold, so about three orders in four (3{,}683 of 5{,}000) include a chilled or frozen
    item and walk to the fridges whatever the plan.
  \item \textbf{The fastest ambient items were already at the front.} The demo layout had eggs (24.9 a day) and bread
    (14.8) in the first two bins; the plan only swaps them.
  \item \textbf{Equal sellers only trade places.} The three chilled items sell at almost the same rate, so curd and
    milk only swap the first two chilled bins. The other ambient items sell 3 to 5 a day, so their order is mostly noise.
\end{itemize}
In this layout the walk is decided by where the fridges stand, not by which shelf holds which item. A bigger saving
needs the fridges nearer the packing desk, which is a layout decision. Slotting should matter more in a real store,
with thousands of SKUs of very different speeds in each zone, than in a demo store with 16.

\section{Trade-off: for sale at putaway, not when the truck is done (HLD 13)}
\label{sec:m3-tradeoff}

The prototype first received a truck in one call: one transaction, all or nothing, and no record of the receipt. That
is the simplest correct design, because if anything fails, nothing happened. HLD 6.3 asks for more, each line for sale
the moment it is in its bin, and task 3.5 built that as five steps with a saved receipt. Building it showed what a
faster shelf costs.

\textbf{What it gains.} A fast seller that has run out is for sale as soon as its case is on the shelf, not after the
last case of the truck is scanned. The receipt also records what was ordered, scanned, damaged, refused, put away and
sent back, which the single call never kept.

\textbf{What it costs.}
\begin{itemize}
  \item \textbf{A third place for stock.} Units on the dock are neither on the shelf nor gone, and every rule had to say
    what happens to them: they still count as in transit, a receipt cannot close while they wait unless they are sent
    back, and a receipt can be cancelled only before anything from it is put away.
  \item \textbf{Races and retries the single call never had.} Two receivers could put away the same line, one could
    close the receipt while another scans, and a handheld could send a scan twice. Each step locks the receipt's row,
    units leave the dock through one conditional \texttt{UPDATE}, a bin's room is checked again after writing, and
    every write takes an \texttt{Idempotency-Key}, so a repeated scan counts once.
  \item \textbf{Edge cases keep appearing.} Review found that a line whose batch number is already in the store with
    another expiry date can be scanned onto the dock but never put away, so its receipt closes only by sending the line
    back. The fix, refusing such a scan at the dock, is waiting for review.
\end{itemize}

We kept the single call as well: \texttt{POST /v1/goods-receipts} runs the five steps in one transaction, so the ops
screen's \emph{Receive truck} button is still all or nothing, and both paths share one set of rules. We accepted a receipt with more states, and more
code to keep them consistent, in exchange for a shelf that refills line by line and a record of every truck.

\section{Limits and open items}

Table~\ref{tab:m3-limits} lists what the supply work leaves open, and what would close each item.

\begin{table}[!htb]
  \centering
  \footnotesize
  \begin{tabular}{@{}>{\raggedright\arraybackslash}p{0.27\linewidth}>{\raggedright\arraybackslash}p{0.36\linewidth}
                  >{\raggedright\arraybackslash}p{0.31\linewidth}@{}}
    \toprule
    Limit & What it means & Next step \\
    \midrule
    Daily forecasts, not hourly & The evening peak is not planned for & Spread each daily forecast over the hours by
      their share of sales \\
    \addlinespace[3pt]
    No festivals or weather & A festival day is missed, then repeated in the same weekday's forecast for four weeks
      & A festival calendar as an input \\
    \addlinespace[3pt]
    Reorder point from today's forecast only & Run on a busy Sunday, milk is sized for a Sunday & Size it from the
      forecasts of the lead-time days \\
    \addlinespace[3pt]
    The same cover for fresh and long-life food & Three days of cover above the reorder point, also for milk that keeps
      three days & Cover days that depend on shelf life \\
    \addlinespace[3pt]
    A batch number already in the store with another expiry date & It is scanned onto the dock but cannot be put away;
      only sending it back closes the receipt & Refuse it at the scan (pull request \#77, pending) \\
    \addlinespace[3pt]
    Slotting only within a zone; apply records the moves at once & 1\% shorter walks in this layout, and staff must
      carry the stock straight away & Fridges nearer the packing desk; items bought together side by side \\
    \bottomrule
  \end{tabular}
  \caption{Limits of the supply work and what would close them.}
  \label{tab:m3-limits}
\end{table}
